\section{Common private source spaces}
\label{sec:peps_common_source_spaces}

\begin{definition}[Sources in prescribed pair slots]
    \label{def:peps_source_slots}
    \lean{TNLean.PEPS.PairEffect.SourceInventory.ofSlots}
    \lean{TNLean.PEPS.PairEffect.SourceInventory.ofSlots_nil}
    \lean{TNLean.PEPS.PairEffect.SourceInventory.ofSlots_cons}
    \lean{TNLean.PEPS.PairEffect.SourceInventory.isNormalized_ofSlots}
    \lean{TNLean.PEPS.PairEffect.SourceInventory.layout_ofSlots_eq}
    \lean{TNLean.PEPS.PairEffect.SourceInventory.partyPairs_ofSlots}
    \lean{TNLean.PEPS.PairEffect.SourceInventory.slotLayout}
    \lean{TNLean.PEPS.PairEffect.SourceInventory.prepareSlots}
    \leanok
    \uses{def:peps_source_inventory}
    Fix an ordered list of pairs $(p_e,q_e)$ with $p_e\ne q_e$.
    Assign complex inner product spaces $U_e,V_e$ to the respective
    parties and a vector $\eta_e\in U_e\otimes V_e$ to each pair.
    These data define an ordered source list. Its register layout depends
    only on the parties and halfspaces, and its unordered pairs are the
    prescribed pairs. The list is normalized when each $\eta_e$ is
    normalized. The empty list prepares no registers; separating the
    first pair separates its two registers from those of the remaining
    sources.

    For any original memory $\mathcal H$, write $Q_{\eta,\mathcal H}$
    for preparation of these sources, using the fixed register layout
    and the canonical associativity identifications. Thus
    $Q_{\eta,\mathcal H}x=(\bigotimes_e\eta_e)\otimes x$.
    These are the pair slots of
    \cite[proof of Theorem~5.2]{OpenAI2026PolynomialPEPS},
    \texttt{04-compression.tex}, lines 253--267.
\end{definition}

\begin{theorem}[A common order and endpoint orientation]
    \label{thm:peps_source_ordering}
    \lean{TNLean.PEPS.PairEffect.SourceInventory.Expands.of_perm}
    \lean{TNLean.PEPS.PairEffect.SourceInventory.exists_ofSlots_expands}
    \leanok
    \uses{def:peps_source_slots,def:peps_source_expansion}
    Let $R$ be a reference source list with distinct unordered pairs,
    and let $S$ be a normalized source list with precisely the same
    distinct unordered pairs. There are normalized sources in the pair
    order and endpoint orientations of $R$ whose preparation can be
    transformed into the preparation of $S$ by allowed operations
    containing no sources, for every spectator memory. Permuting whole
    source blocks also gives such a transformation.
\end{theorem}

\begin{proof}\leanok
    \uses{def:peps_source_expansion}
    Induct on the reference list. Select the actual source on its first
    unordered pair and move this source block to the front. If its
    endpoint order is reversed, apply the tensor-factor exchange
    $U\otimes V\to V\otimes U$, which preserves the vector norm.
    Remove the matched pair and continue. Equal pair sets and
    distinctness ensure that every source is used exactly once.
\end{proof}

\begin{definition}[Orthogonal sums of branch spaces]
    \label{def:peps_common_halfspaces}
    \lean{TNLean.PEPS.PairEffect.PairSource.commonSpace}
    \lean{TNLean.PEPS.PairEffect.PairSource.commonInclusion}
    \lean{TNLean.PEPS.PairEffect.PairSource.commonProjection}
    \leanok
    Let $\Xi$ be a finite set and let $U_\xi$ be complex inner product
    spaces. Their common space is the orthogonal sum
    $\widehat U=\bigoplus_{\xi\in\Xi}U_\xi$.
    Denote its coordinate inclusion by
    $j_\xi:U_\xi\to\widehat U$ and its coordinate projection by
    $\pi_\xi:\widehat U\to U_\xi$.
    The norm is given by
    $\|(u_\xi)_{\xi\in\Xi}\|^2=\sum_{\xi\in\Xi}\|u_\xi\|^2$.
    This is the common-space construction in
    \cite[proof of Theorem~5.2]{OpenAI2026PolynomialPEPS},
    \texttt{04-compression.tex}, lines 253--267.
\end{definition}

\begin{theorem}[Embedding and recovering a pair source]
    \label{thm:peps_common_pair_source}
    \lean{TNLean.PEPS.PairEffect.PairSource.commonProjection_commonInclusion}
    \lean{TNLean.PEPS.PairEffect.PairSource.norm_commonProjection_le}
    \lean{TNLean.PEPS.PairEffect.PairSource.commonVector}
    \lean{TNLean.PEPS.PairEffect.PairSource.norm_commonVector}
    \lean{TNLean.PEPS.PairEffect.PairSource.mapL_commonProjection_commonVector}
    \leanok
    \uses{def:peps_common_halfspaces}
    Each $j_\xi$ is isometric, $\|\pi_\xi\|\le1$, and
    $\pi_\xi j_\xi=\Id$. For a second family $V_\xi$, use inclusions
    $k_\xi$ and projections $\rho_\xi$ on its orthogonal sum.
    If $\eta_\xi\in U_\xi\otimes V_\xi$, set
    $\widehat\eta_\xi=(j_\xi\otimes k_\xi)\eta_\xi$.
    Then
    \begin{align}
        \|\widehat\eta_\xi\|                     & =\|\eta_\xi\|, \\
        (\pi_\xi\otimes\rho_\xi)\widehat\eta_\xi & =\eta_\xi.
        \label{eq:peps_common_pair_recovery}
    \end{align}
    In particular, normalized sources remain normalized in the common
    spaces.
\end{theorem}

\begin{proof}\leanok
    A coordinate inclusion has exactly one nonzero coordinate, so it
    preserves norms. Every coordinate norm is bounded by the norm of
    the orthogonal sum, proving the projection bound. Tensor products
    of linear isometries preserve norms. The recovery identity follows
    first on elementary tensors from $\pi_\xi j_\xi=\Id$ and
    $\rho_\xi k_\xi=\Id$, then by linearity.
\end{proof}

\begin{theorem}[Local maps on a prepared pair]
    \label{thm:peps_map_source_pair}
    \lean{TNLean.PEPS.PairEffect.Word.mapPair}
    \lean{TNLean.PEPS.PairEffect.Word.mapPair_spec}
    \lean{TNLean.PEPS.PairEffect.Word.eval_mapPair_tmul}
    \lean{TNLean.PEPS.PairEffect.Word.eval_mapPair_source}
    \leanok
    \uses{def:peps_party_layout,def:peps_source_inventory}
    Let $A:U\to U'$ and $B:V\to V'$ be continuous linear maps at the
    respective parties of a pair source. Applying these maps locally,
    with identities on a spectator memory, sends
    $u\otimes(v\otimes x)$ to $Au\otimes(Bv\otimes x)$ and satisfies
    \begin{align}
        (A\otimes B\otimes\Id)P_{\eta,\mathcal H}
         & =P_{(A\otimes B)\eta,\mathcal H}.
        \label{eq:peps_map_source_pair}
    \end{align}
    When $A$ and $B$ are contractions, these two local operations form
    an allowed composition containing no sources.
\end{theorem}

\begin{proof}\leanok
    Act first at the first party and then at the second party, keeping
    the remaining registers untouched. The formula follows on
    elementary tensors, hence on every source vector by linearity.
    Tensoring each local contraction with an identity preserves its
    contraction bound.
\end{proof}

\begin{theorem}[Common spaces for normalized source lists]
    \label{thm:peps_common_source_lists}
    \lean{TNLean.PEPS.PairEffect.SourceInventory.common_ofSlots_expands}
    \lean{TNLean.PEPS.PairEffect.SourceInventory.exists_common_expansions}
    \lean{TNLean.PEPS.PairEffect.SourceInventory.exists_prepareSlots_recovery}
    \lean{TNLean.PEPS.PairEffect.SourceInventory.exists_prepareSlots_expands}
    \leanok
    \uses{def:peps_source_slots,def:peps_common_halfspaces,def:peps_source_expansion}
    Let $(S_\xi)_{\xi\in\Xi}$ be a finite family of normalized source
    lists, each with the same distinct unordered pair set as a fixed
    reference list. There are common halfspaces in the reference pair
    order and normalized vectors $\widehat\eta_{\xi,e}$ such that,
    for every spectator memory $\mathcal H$, an allowed composition
    $D_\xi$ containing no sources satisfies
    \begin{align}
        D_\xi Q_{\widehat\eta_\xi,\mathcal H}
         & =P_{S_\xi,\mathcal H}.
        \label{eq:peps_common_list_recovery}
    \end{align}
    The domain of $D_\xi$ is the fixed source register layout together
    with $\mathcal H$; it does not depend on the vectors.
\end{theorem}

\begin{proof}\leanok
    \uses{thm:peps_source_ordering,thm:peps_common_pair_source,thm:peps_map_source_pair}
    Put every list in the reference order using
    Theorem~\ref{thm:peps_source_ordering}. At each half of each pair,
    form the orthogonal sum of its branch spaces and embed the branch
    vector. Apply the coordinate projections to each pair in turn.
    Equations~(\ref{eq:peps_common_pair_recovery}) and
    (\ref{eq:peps_map_source_pair}) recover its original source vector.
    Compose these operations with the exchanges recovering the original
    source order. All local maps remain contractions.
\end{proof}

\begin{theorem}[Local factorization in common private spaces]
    \label{thm:peps_common_party_maps}
    \lean{TNLean.PEPS.PairEffect.Word.exists_common_source_preparation}
    \lean{TNLean.PEPS.PairEffect.Word.exists_common_prepared_tensorPartyMaps}
    \leanok
    \uses{def:peps_source_slots,def:peps_party_registers}
    Let $P$ be the prescribed finite set of gate parties, with a fixed
    ordering, and let $(F_\xi)_{\xi\in\Xi}$ be a finite family of
    allowed compositions with the same input and output register
    layouts. There is one common source slot for every unordered pair
    of distinct parties, with fixed halfspaces, normalized vectors
    $\widehat\eta_{\xi,e}$, and local contractions $B_{\xi,p}$ such that
    \begin{align}
        U_{\rm out}F_\xi
         & =\left(\bigotimes_{p\in P}B_{\xi,p}\right)
        U_{\rm in}Q_{\widehat\eta_\xi,\mathcal H_{\rm in}}.
        \label{eq:peps_common_party_maps}
    \end{align}
    Here $U_{\rm in}$ groups the common source registers and the original
    input registers by party, and $U_{\rm out}$ groups the original
    output registers in the same party order. Both identifications and
    the domain and codomain of each party's local maps are independent
    of $\xi$. Equivalently, preparation in the common source layout can
    be followed by an allowed composition containing no sources to
    recover $F_\xi$ exactly.

    These are the common private spaces in
    \cite[proof of Theorem~5.2]{OpenAI2026PolynomialPEPS},
    \texttt{04-compression.tex}, lines 253--267.
\end{theorem}

\begin{proof}\leanok
    \uses{thm:peps_complete_pair_sources,thm:peps_complete_prepared_party_maps,
        thm:peps_common_source_lists,thm:peps_party_factorization}
    Choose a reference ordering of all unordered pairs by completing
    the empty source list. Complete the actual sources of each $F_\xi$
    by Theorem~\ref{thm:peps_complete_prepared_party_maps}, and then apply
    Theorem~\ref{thm:peps_common_source_lists}. Compose the resulting
    recovery maps with the remaining operations of $F_\xi$ and collect
    the local maps by party. This gives (\ref{eq:peps_common_party_maps}).
    The reference list is chosen independently of $\xi$, so the
    construction also applies when $\Xi$ is empty.
\end{proof}

\begin{theorem}[The weighted gate in common private spaces]
    \label{thm:peps_common_source_gate}
    \lean{TNLean.PEPS.PairEffect.Word.exists_common_source_gate}
    \leanok
    \uses{def:peps_source_slots,def:peps_party_registers}
    Let $F=\sum_{\xi\in\Xi}c_\xi F_\xi$, where $\Xi$ is finite,
    each $c_\xi\in\C$, and the $F_\xi$ are allowed compositions with
    the same input and output register layouts on a prescribed finite
    set of parties $P$. There are fixed halfspaces for all unordered
    pairs of distinct parties, normalized pair sources
    $\widehat\eta_{\xi,e}$, and local contractions $B_{\xi,p}$ such that
    \begin{align}
        U_{\rm out}F
         & =\sum_{\xi\in\Xi}c_\xi
        \left(\bigotimes_{p\in P}B_{\xi,p}\right)
        U_{\rm in}Q_{\widehat\eta_\xi,\mathcal H_{\rm in}}.
        \label{eq:peps_common_source_gate}
    \end{align}
    The coefficients are those of the original sum. The input and
    output identifications and all private spaces are common to every
    summand. This is the common-space form of
    \cite[proof of Theorem~5.2]{OpenAI2026PolynomialPEPS},
    \texttt{04-compression.tex}, lines 233--267.
\end{theorem}

\begin{proof}\leanok
    \uses{thm:peps_common_party_maps}
    Apply Theorem~\ref{thm:peps_common_party_maps} to the family
    $(F_\xi)_{\xi\in\Xi}$. Multiply each identity
    (\ref{eq:peps_common_party_maps}) by $c_\xi$ and sum over $\Xi$.
\end{proof}


\begin{theorem}[Finite coordinates for a finite family of source vectors]
    \label{thm:peps_finite_pair_coordinates}
    \lean{TNLean.PEPS.PairEffect.PairSource.exists_finite_coordinates}
    \leanok
    Let $U,V$ be complex inner product spaces, let $\Xi$ be finite, and
    let $\eta_\xi\in U\otimes V$ for each $\xi\in\Xi$, where the tensor
    product is algebraic. There exist $a,b\in\N$, linear isometries
    $f:\C^a\to U$ and $g:\C^b\to V$, and vectors
    $\eta^0_\xi\in\C^a\otimes\C^b$ such that
    $(f\otimes g)\eta^0_\xi=\eta_\xi$ and
    $\|\eta^0_\xi\|=\|\eta_\xi\|$ for every $\xi$.
    The integers $a,b$ may be zero. This gives finite coordinates for
    the source vectors used in the Schmidt decompositions of
    \cite[proof of Theorem~5.2]{OpenAI2026PolynomialPEPS},
    \texttt{04-compression.tex}, lines 279--299.
\end{theorem}

\begin{proof}\leanok
    A finite set of algebraic tensors is contained in $U_0\otimes V_0$
    for some finite-dimensional subspaces $U_0\subseteq U$ and
    $V_0\subseteq V$: express each tensor as a finite sum of elementary
    tensors and span the finitely many factors on each side.
    Choose orthonormal coordinates on $U_0,V_0$ and compose their
    coordinate isometries with the subspace inclusions. The induced
    tensor map is isometric, so the coordinate vectors have exactly
    the original norms.
\end{proof}

\begin{theorem}[Source preparations in finite coordinates]
    \label{thm:peps_finite_source_coordinates}
    \lean{TNLean.PEPS.PairEffect.SourceInventory.ofSlots_mapIsometry_expands}
    \lean{TNLean.PEPS.PairEffect.SourceInventory.exists_finite_coordinate_expansions}
    \leanok
    \uses{def:peps_source_slots,def:peps_source_expansion}
    Fix a finite source list, common halfspaces $U_e,V_e$, and a finite
    family of vectors $\eta_{\xi,e}\in U_e\otimes V_e$.
    There are dimensions $a_e,b_e$, isometries
    $f_e:\C^{a_e}\to U_e$ and $g_e:\C^{b_e}\to V_e$, and coordinate
    vectors $\eta^0_{\xi,e}$ such that
    $(f_e\otimes g_e)\eta^0_{\xi,e}=\eta_{\xi,e}$ and
    $\|\eta^0_{\xi,e}\|=\|\eta_{\xi,e}\|$.
    For every $\xi$ and every spectator memory $\mathcal H$, an allowed
    composition $D_\xi$ containing no sources satisfies
    $D_\xi Q_{\eta^0_\xi,\mathcal H}=Q_{\eta_\xi,\mathcal H}$. More generally, any prescribed isometric inclusions of the
    two halves of each source induce such a recovery of the included
    sources.
\end{theorem}

\begin{proof}\leanok
    \uses{thm:peps_finite_pair_coordinates,thm:peps_map_source_pair}
    Apply Theorem~\ref{thm:peps_finite_pair_coordinates} at each pair
    slot. Apply $f_e$ and $g_e$ locally to the two halves of that slot,
    then proceed through the source list. Equation
    (\ref{eq:peps_map_source_pair}) recovers the corresponding source
    vector at each step. The inclusions are isometries, hence
    contractions, and create no further sources.
\end{proof}

\begin{theorem}[The source-gate expansion in finite coordinates]
    \label{thm:peps_finite_source_gate}
    \lean{TNLean.PEPS.PairEffect.Word.exists_finite_source_preparation}
    \lean{TNLean.PEPS.PairEffect.Word.exists_finite_prepared_tensorPartyMaps}
    \lean{TNLean.PEPS.PairEffect.Word.exists_finite_source_gate}
    \leanok
    \uses{def:peps_source_slots,def:peps_party_registers}
    Let $F=\sum_{\xi\in\Xi}c_\xi F_\xi$, where $\Xi$ is finite,
    $c_\xi\in\C$, and each $F_\xi$ is an allowed composition between
    fixed register layouts on a prescribed finite set of parties $P$.
    There are common dimensions $a_e,b_e$ for every unordered pair of
    distinct parties, normalized sources
    $\eta^0_{\xi,e}\in\C^{a_e}\otimes\C^{b_e}$, and local
    contractions $B_{\xi,p}$ such that
    \begin{align}
        U_{\rm out}F_\xi
         & =\left(\bigotimes_{p\in P}B_{\xi,p}\right)
        U_{\rm in}Q_{\eta^0_\xi,\mathcal H_{\rm in}}, \\
        U_{\rm out}F
         & =\sum_{\xi\in\Xi}c_\xi
        \left(\bigotimes_{p\in P}B_{\xi,p}\right)
        U_{\rm in}Q_{\eta^0_\xi,\mathcal H_{\rm in}}.
        \label{eq:peps_finite_source_gate}
    \end{align}
    Here $U_{\rm in}$ groups the common finite source registers and
    original input registers by party, while $U_{\rm out}$ groups the
    unchanged output registers in the same order. Each branch is also
    the exact composition of its coordinate source preparation with
    allowed operations containing no sources.

    This is an exact finite-coordinate form of
    \cite[proof of Theorem~5.2]{OpenAI2026PolynomialPEPS},
    \texttt{04-compression.tex}, lines 233--299. No bound on the private
    dimensions $a_e,b_e$ is asserted.
\end{theorem}

\begin{proof}\leanok
    \uses{thm:peps_common_party_maps,thm:peps_complete_prepared_party_maps,
        thm:peps_common_source_lists,thm:peps_finite_source_coordinates,
        thm:peps_party_factorization}
    Theorem~\ref{thm:peps_common_party_maps} prepares every branch's
    sources in common private spaces, followed by allowed operations
    containing no sources. Apply
    Theorem~\ref{thm:peps_finite_source_coordinates} to the finitely many
    common vectors at each slot. Compose the coordinate inclusions with
    the remaining operations of each branch. All these operations are local
    contractions. Collect them by party and then sum the resulting
    branch identities with the original coefficients $c_\xi$.
\end{proof}
